\documentclass[UTF8,zihao=-4]{ctexart}

\usepackage[a4paper,margin=2.3cm]{geometry}
\usepackage{array}
\usepackage{booktabs}
\usepackage{longtable}
\usepackage{hyperref}
\usepackage{xcolor}
\usepackage{enumitem}
\usepackage{fancyhdr}
\usepackage{lastpage}

\hypersetup{
  colorlinks=true,
  linkcolor=black,
  urlcolor=blue
}

\setlength{\parindent}{2em}
\setlength{\parskip}{0.35em}
\renewcommand{\arraystretch}{1.25}
\setlist[itemize]{leftmargin=2.2em,itemsep=0.2em,topsep=0.2em}
\pagestyle{fancy}
\fancyhf{}
\fancyfoot[C]{第 \thepage 页 / 共 \pageref*{LastPage} 页}
\renewcommand{\headrulewidth}{0pt}
\renewcommand{\footrulewidth}{0pt}

\title{\heiti 托福阅读 Class 4 课堂笔记}
\author{}
\date{}

\begin{document}
\maketitle
\thispagestyle{fancy}

本节课前半部分整理托福阅读中的一般生词、易混词和同义替换，后半部分系统梳理与“石头”“超越”“时间”“名称”以及“古老、原始、最初”相关的词根词缀。整理时保留课堂上的核心拆解、常考词义和高频搭配，方便在阅读中快速识别陌生词并判断其语境含义。

\section*{一、一般生词、易混词与同义替换}

\texttt{scale} 既可表示“鱼鳞”，也可表示“规模”，托福阅读中常见搭配为 \texttt{on a large scale}；\texttt{meager} 表示“微薄的、贫乏的”；\texttt{weathering} 是地质学文章中的高频词，表示“风化（作用）”；\texttt{ingenious} 表示“有创造力的、精巧的”，其名词 \texttt{ingenuity} 指“独创性、聪明才智”。

形近词中，\texttt{indigenous} 表示“本土的、土著的”，\texttt{indigent} 表示“贫困的”，\texttt{igneous rock} 表示“火成岩”。这组三词在拼写上接近，但语义领域完全不同，阅读时尤其需要结合上下文迅速区分。与“繁荣、生长”相关的同义替换中，\texttt{flourish}、\texttt{thrive} 与 \texttt{prosper} 都可表示“繁荣、茂盛、兴旺”；\texttt{fecund} 则强调“肥沃、多产、富于创造力”，在农业、生态和文明发展类文章里常可替换 \texttt{fertile} 或 \texttt{productive}。

\texttt{sustain} 表示“维持、支撑”，其名词 \texttt{sustenance} 表示“食物、营养、维持生命之物”；\texttt{sophistication} 可表示“复杂、精密、老练”。\texttt{subjugate} 表示“征服、使服从”，可拆为 \texttt{sub-} “在……之下” + 与“轭、束缚”相关的词根 + 动词后缀，核心意象是“把别人置于控制之下”。\texttt{mold} 是典型的一词多义词，既可表示“模具”，也可表示“霉菌”，作动词时还有“塑造”之意。

另一组常见形近词是 \texttt{industrial} 与 \texttt{industrious}：前者表示“工业的、产业的”，后者表示“勤奋的、勤勉的”。短语 \texttt{be long overdue in doing} 表示“早就该做某事而一直拖到现在”，其中 \texttt{overdue} 本义是“过期的、延误的”，在阅读和写作中都很实用。

\section*{二、与“石头”和“超越”相关的词根词缀}

\texttt{-lith-} 与“石头”有关，因此 \texttt{the Neolithic Age}、\texttt{the Mesolithic Age} 与 \texttt{the Paleolithic Age} 分别表示“新石器时代”“中石器时代”和“旧石器时代”；\texttt{lithiasis} 表示“结石”；\texttt{monolith} 表示“独石柱”或“单一而庞大的整体”；\texttt{lithography} 表示“石版印刷术”；\texttt{lithophyte} 表示“岩生植物”。课堂还顺带扩展了一组以 \texttt{-sphere} 结尾的地球科学词汇，如 \texttt{lithosphere} “岩石圈”、\texttt{hydrosphere} “水圈”、\texttt{biosphere} “生物圈”、\texttt{atmosphere} “大气层”、\texttt{troposphere} “对流层”、\texttt{stratosphere} “平流层”、\texttt{mesosphere} “中间层”、\texttt{thermosphere} “热层”和 \texttt{exosphere} “外逸层”。

前缀 \texttt{sur-} 常表示“在上方、超过、越过”。\texttt{surrealism} 是“超现实主义”，\texttt{surcharge} 是“附加费”，\texttt{surtax} 是“附加税”，\texttt{surmount} 表示“克服、越过”，因此 \texttt{insurmountable} 表示“难以克服的”；\texttt{surrender} 表示“投降”；\texttt{survive} 则可理解为“活过、活得比……更久”，即“幸存”。

表示“超出、极限、过度”的前缀群包括 \texttt{extra-}、\texttt{ultra-}、\texttt{hyper-} 与 \texttt{super-}。\texttt{extra-} 常表示“在……之外、超出”，如 \texttt{extraordinary} “非凡的”、\texttt{extraterrestrial} “地球外的、外星的”、\texttt{extravagant} “奢侈的、铺张的”、\texttt{extracellular} “细胞外的”；\texttt{ultra-} 强调“极度、超出”，如 \texttt{ultraviolet} “紫外线”、\texttt{ultrasonic} “超声波的”、\texttt{ultralight} “超轻的”；\texttt{hyper-} 常表示“过度、异常高”，如 \texttt{hyperlink}、\texttt{hypertext}、\texttt{hyperbole}、\texttt{hypertension}、\texttt{hyperglycemia}、\texttt{hyperlipidemia}、\texttt{hypercholesterolemia} 与 \texttt{hyperthyroidism}；\texttt{super-} 则常表示“在上面、在更高层级”，如 \texttt{superimpose} “叠加”和 \texttt{superconductivity} “超导性”。

\section*{三、与“时间”“名称”相关的核心词根}

\texttt{-chron-} 表示“时间”。\texttt{chronicle} 是“编年史、记录”，\texttt{chronic} 是“慢性的、长期的”，\texttt{chronobiology} 是“时间生物学”，\texttt{chronograph} 是“秒表、精密计时器”。\texttt{synchronize} 及其名词 \texttt{synchronization} 表示“同步”，\texttt{desynchronization} 表示“去同步化、失步”，\texttt{chronological} 表示“按时间顺序的”，\texttt{chronemics} 是“时间行为学”，\texttt{chronophage} 则比喻“吞噬时间的事物”。课堂还补充了 \texttt{monochronic} 与 \texttt{polychronic} 两种时间管理方式，以及 \texttt{live in sync with the time}、\texttt{out of sync with} 等固定表达。

\texttt{-nym- / -nom-} 源自希腊语，表示“名字、词名”。\texttt{synonymous} 表示“同义的、等同于……的”，常见搭配是 \texttt{be synonymous with}；\texttt{synonym} 与 \texttt{antonym} 分别表示“同义词”和“反义词”；\texttt{pseudonym} 是“笔名、假名”；\texttt{acronym} 是“首字母缩略词”；\texttt{heteronym} 指“同形异音异义词”，\texttt{homonym} 指“同音异义词”；\texttt{eponym} 指“词源人名、同名命名现象”，课堂用 \texttt{Xerox} 由品牌名延伸为“复印”动作作了说明。另一方面，\texttt{anonymous} 表示“匿名的、无名的”，对应名词为 \texttt{anonymity}；\texttt{nominal} 表示“名义上的”或“数额很小的”；\texttt{nominate}、\texttt{nomination} 与 \texttt{nominee} 则对应“提名、提名行为、被提名者”。

\section*{四、表示“古老、原始、最初”的前缀群与课堂总结}

前缀 \texttt{arch- / archaeo-} 常表示“古老、原始、首要”。\texttt{archaeology} 是“考古学”，\texttt{archaic} 表示“古老的、过时的”，对应名词 \texttt{archaism} 为“古语、古体”；\texttt{archive} 现在表示“档案馆、档案”或“存档”，词源上最初与“首长办公室”有关；\texttt{archetype} 表示“原型、典型”，强调最初的模型或范式。

\texttt{paleo-} 表示“远古的、史前的”，因此 \texttt{paleontology} 是“古生物学”，\texttt{paleoclimate} 是“古气候”，\texttt{paleoanthropology} 是“古人类学”，\texttt{Paleozoic} 是“古生代”，而 \texttt{the Paleolithic Age} 则回扣到了前文讲到的“旧石器时代”。\texttt{proto-} 表示“第一、最原始、雏形”，因此 \texttt{prototype} 是“原型、雏形、最初型号”，\texttt{protozoon / protozoa} 表示“原生动物”。这一组前缀在考古学、地质学、生物学和人类学阅读中都非常高频，掌握后能帮助快速判断文章讨论的是“远古阶段”“原始形态”还是“最初模型”。

\section*{五、补充词汇与固定表达}

\begin{longtable}{>{\raggedright\arraybackslash}p{4.9cm} >{\raggedright\arraybackslash}p{9.6cm}}
\toprule
\textbf{单词或表达} & \textbf{含义或用法} \\
\midrule
\endfirsthead
\toprule
\textbf{单词或表达} & \textbf{含义或用法} \\
\midrule
\endhead
\texttt{on a large scale} & 大规模地；从大范围上看 \\
\texttt{meager} & 微薄的，贫乏的 \\
\texttt{weathering} & 风化（作用） \\
\texttt{ingenious} & 有创造力的，精巧的 \\
\texttt{ingenuity} & 独创性，聪明才智 \\
\texttt{indigenous} & 本土的，土著的 \\
\texttt{indigent} & 贫困的 \\
\texttt{igneous rock} & 火成岩 \\
\texttt{flourish / thrive / prosper} & 繁荣，兴旺，茁壮成长 \\
\texttt{sustenance} & 食物，营养，维持生命之物 \\
\texttt{fecund} & 肥沃的，多产的，富于创造力的 \\
\texttt{subjugate} & 征服，使服从 \\
\texttt{mold} & 模具；霉菌；塑造 \\
\texttt{industrial / industrious} & 工业的、产业的 / 勤奋的、勤勉的 \\
\texttt{be long overdue in doing} & 早就该做某事而拖到现在 \\
\texttt{insurmountable} & 难以克服的 \\
\texttt{lithosphere} & 岩石圈 \\
\texttt{extraordinary} & 非凡的，不寻常的 \\
\texttt{extraterrestrial} & 地球外的，外星的 \\
\texttt{ultraviolet} & 紫外线 \\
\texttt{hypertension} & 高血压 \\
\texttt{chronological} & 按时间顺序的 \\
\texttt{chronophage} & 吞噬时间的事物，极耗时间的事物 \\
\texttt{be synonymous with} & 与……同义；与……等同 \\
\texttt{pseudonym} & 笔名，假名 \\
\texttt{acronym} & 首字母缩略词 \\
\texttt{anonymous / anonymity} & 匿名的、无名的 / 匿名状态 \\
\texttt{nominal} & 名义上的；数额很小的 \\
\texttt{archetype} & 原型，典型 \\
\texttt{paleontology} & 古生物学 \\
\texttt{prototype} & 原型，雏形，最初型号 \\
\texttt{protozoon / protozoa} & 原生动物（单数 / 复数） \\
\bottomrule
\end{longtable}

\end{document}
